\newpage
\section{Résultats et Analyse}
\label{sec:benchmark}

Cette section rend compte de l'état d'avancement de la tâche. Elle rapporte les résultats effectivement établis, qui portent sur les briques dont dépend le protocole de la section~\ref{sec:methodo}, et délimite ensuite ce qui ne l'est pas. Cette séparation est explicite parce que l'enjeu de la tâche 3.2 est la fiabilité des explications : un rapport qui présenterait comme acquis un classement de méthodes non mesuré illustrerait exactement le défaut qu'il prétend traiter.

\subsection{Validation du socle de données}
\label{sec:validation-indices}

Le protocole de la section~\ref{sec:methodo} repose sur des indicateurs standardisés, employés à la fois comme variables d'entrée possibles et comme formulation de la cible contrefactuelle. Leur justesse conditionne donc celle de tout ce qui s'appuie dessus, et elle a été vérifiée.

Un indice standardisé doit, par construction, suivre approximativement une loi normale centrée réduite sur sa propre période de référence. Le critère d'acceptation a été fixé \emph{avant} la mesure : moyenne voisine de zéro à $\pm\,\num{0,05}$, écart-type voisin de un à $\pm\,\num{0,05}$, et taux de saturation aux bornes inférieur à \SI{1}{\percent}. La mesure porte sur \num{41960400} valeurs couvrant \num{11496} mailles.

\begin{table}[H]
\centering
\caption{Vérification des indices produits sur leur propre période de référence, \num{41960400} valeurs couvrant \num{11496} mailles. Le critère d'acceptation avait été fixé avant la mesure.}
\label{tab:validation-indices}
\begin{tabularx}{\textwidth}{@{}lXrrrr@{}}
\toprule
\textbf{Indice} & \textbf{Fenêtres} & \textbf{Moyenne} & \textbf{Écart-type} & \textbf{Médiane}
& \textbf{Saturation} \\
\midrule
SPI  & 1/3/6/12 & $-0{,}008$ à $+0{,}002$ & \numrange{0,985}{1,031} & \numrange{0,00}{0,04} & \SIrange{0,05}{0,51}{\percent} \\
STI  & 1/3/6/12 & $0{,}000$               & \num{0,983}            & $-0{,}09$ à $-0{,}02$ & \SIrange{0,00}{0,04}{\percent} \\
SPEI & 1/3/6/12 & $+0{,}004$ à $+0{,}006$ & \numrange{0,999}{1,013} & \numrange{0,01}{0,03} & \SIrange{0,012}{0,035}{\percent} \\
\bottomrule
\end{tabularx}
\end{table}

Les trois indices satisfont le critère sur les quatre fenêtres d'agrégation, et la couverture du SPEI est uniforme à \SI{100}{\percent} après le changement de loi rapporté en section~\ref{sec:arbitrages}, ce qui signifie que la calibration n'est pas seulement préservée mais mesurée sur l'ensemble du domaine.

Deux lectures de ces indices doivent être écartées, et elles sont consignées ici parce que l'erreur correspondante est naturelle. La classe \emph{normale} ne couvre pas la proportion théorique : elle représente \SI{54,8}{\percent} des valeurs au lieu des \SI{59,9}{\percent} attendus d'une gaussienne, alors même que l'écart-type vaut \num{1,008}, la distribution conservant des épaules légèrement plus lourdes. Cet écart est une propriété de la transformation au voisinage des bornes de classe, sans conséquence opérationnelle, et il ne doit pas être lu comme une anomalie de sécheresse. La décroissance décennale du SPEI, en revanche, est un signal climatique et non une dérive de méthode : la moyenne sur fenêtre d'un mois décroît de manière monotone, $+\num{0,038}$ dans les années 1990, $+\num{0,007}$ dans les années 2000, $-\num{0,016}$ dans les années 2010, $-\num{0,117}$ depuis 2020. Cette évolution avait été anticipée avant d'être mesurée, ce qui la distingue d'un artefact d'ajustement.

\subsection{Une incohérence de forçage qui limite l'interprétation du référent}

Le dispositif de la section~\ref{sec:methodo} fait du modèle mécanistique le référent contre lequel juger les explications produites sur un modèle appris. Une incohérence de la chaîne de données en limite aujourd'hui la portée, et elle doit être rapportée ici plutôt que dissimulée dans les annexes techniques, parce qu'elle porte directement sur la validité du protocole.

Le forçage d'évapotranspiration transporté par la chaîne des stations est la variable brute d'ERA5-Land et non l'évapotranspiration de référence, soit environ le double de la grandeur attendue, \SI{1839}{\milli\metre\per\an} mesuré en base contre \SI{818}{\milli\metre\per\an} pour la référence. La correction n'est pas un changement de colonne : la chaîne station ne transporte que la température moyenne, alors que la formule de Hargreaves exige aussi les extrêmes journaliers, de sorte qu'il faudrait les y ajouter, recalculer l'évapotranspiration au pas journalier, propager le résultat dans deux tables partitionnées et recalibrer les modèles existants. Ce que cet écart change pour les calibrations n'a pas été évalué : aucune campagne comparant les deux forçages sur un échantillon de stations n'a été conduite, et l'ampleur de l'effet reste donc inconnue.

La conséquence pour la tâche 3.2 est directe. Les calibrations mécanistiques restent utilisables pour la prévision et pour le test de cohérence croisée, qui compare deux modèles entre eux et non un paramètre à sa valeur physique. En revanche, les paramètres de ces modèles, temps de réponse et gain, ne s'interprètent pas physiquement tant que ce forçage n'est pas corrigé. Toute lecture hydrogéologique d'une fonction de réponse calibrée, et donc toute explication qui s'appuierait sur elle pour justifier une prévision, est suspendue à cette correction. C'est le premier chantier que ces travaux laissent à leur suite, et il est prioritaire sur toute campagne d'évaluation.

\subsection{État de l'outillage au regard des deux questions de recherche}

Le tableau~\ref{tab:etat-qr} récapitule ce que le socle apporte à chaque question et ce qui manque encore. La colonne d'état est délibérément sévère : \emph{outillé} signifie que les moyens existent, pas qu'un résultat a été produit, et la mention \emph{exposé} distingue ce qu'un utilisateur peut atteindre depuis l'application de ce qui ne s'exécute que par l'API ou un carnet de calcul.

\begin{table}[H]
\centering
\caption{État d'avancement au regard des questions de recherche et des questions techniques.}
\label{tab:etat-qr}
\begin{tabularx}{\textwidth}{@{}L{2.2cm}YL{3.1cm}@{}}
\toprule
\textbf{Question} & \textbf{Ce qui a été produit} & \textbf{État} \\
\midrule
QT1 & Entrepôt unifié, quatre sources jointes, indicateurs standardisés validés, chaîne
idempotente et reprenable, déploiement conteneurisé & atteint \\
QT2 & Observatoire, laboratoire mécanistique, catalogue d'architectures profondes, traçabilité
des expériences, modules d'explication & atteint \\
QR1 & Cinq familles d'attribution mises en œuvre et confrontables entre elles, restitution sur
le plan variable-temps ; quatre d'entre elles accessibles dans l'interface & outillé, exposé,
non évalué \\
QR2 & Deux méthodes contrefactuelles de natures différentes, test de cohérence croisée contre
un modèle indépendant, protocole de comparaison formalisé & outillé, non exposé dans
l'interface, non mesuré \\
\bottomrule
\end{tabularx}
\end{table}

\subsection{Ce qui n'est pas établi}

Quatre choses que le lecteur pourrait attendre ne figurent pas dans ce rapport, et il faut dire pourquoi.

Le \textbf{pilotage des contrefactuels depuis l'interface} n'est pas livré. La page correspondante existe dans le dépôt et l'API répond, mais aucune route de l'application n'y mène, de sorte que les méthodes de la section~\ref{sec:xai} restent hors de portée d'un utilisateur non développeur. C'est un défaut de branchement, court à corriger, et non une lacune de méthode ; il est signalé ici parce qu'il change ce que la plateforme permet aujourd'hui de faire, et parce qu'il conditionne toute évaluation associant des hydrogéologues.

Le \textbf{classement des architectures de prévision} sur le cas hydrogéologique n'est pas établi. Les campagnes d'entraînement conduites au cours des travaux ont servi à valider le dispositif et à orienter les développements ; les jeux de données préparés et les points de sauvegarde ont été purgés des dépôts lors de la passation. Reconstituer a posteriori un tableau de performances à partir de souvenirs ne serait pas un résultat scientifique. La campagne qui opposerait les modèles appris aux modèles mécanistiques sur un échantillon stratifié reste donc à conduire, et l'outillage nécessaire est en place pour la mener.

La \textbf{comparaison quantitative des familles de contrefactuels} n'est pas davantage établie. Les deux méthodes décrites en section~\ref{sec:xai} sont opérationnelles et mettent en œuvre l'exigence de plausibilité formulée par les hydrogéologues, mais aucune évaluation n'a été conduite, sur aucun échantillon, selon aucune des dimensions du tableau~\ref{tab:dimensions}. Le protocole de la section~\ref{sec:methodo} est le livrable de cette partie de la tâche : il fixe les conventions qui manquaient, il n'en rapporte pas encore l'application. Anticiper son résultat reviendrait à préjuger de ce que la mesure doit trancher.

L'\textbf{évaluation de la restitution auprès des experts}, enfin, qui constitue la seconde moitié de QR1, n'a pas été conduite. Les visualisations existent et ont été construites en dialogue avec les hydrogéologues du BRGM, mais aucun protocole d'évaluation avec utilisateurs n'a été mis en place. C'est une question distincte des précédentes, qui relève de méthodes d'évaluation propres et qui ne se déduit d'aucune métrique calculable sur les explications elles-mêmes.

\subsection{Analyse}

Le principal enseignement de cette première phase n'est pas d'ordre algorithmique. Il tient à ce que le cas hydrogéologique a révélé des exigences que la littérature XAI sur séries temporelles ne traite qu'incidemment.

La première est que la \textbf{plausibilité ne se pondère pas}. La littérature en fait un terme supplémentaire de la fonction objectif, arbitré contre la validité et la proximité ; sur des variables liées par des relations physiques, un tel arbitrage produit des scénarios qui violent ces relations d'autant moins que le poids est élevé, mais qui les violent toujours. Or un expert rejette un scénario qui les viole un peu aussi sûrement qu'un scénario qui les viole beaucoup : ce n'est pas une grandeur continue pour lui, c'est un test binaire de recevabilité. La conséquence de conception est celle décrite en section~\ref{sec:xai}, changer d'espace de recherche au lieu d'ajuster une pondération. Sa conséquence expérimentale est une question ouverte, que l'évaluation devra trancher : cette garantie s'achète au prix d'une expressivité réduite à sept degrés de liberté, et rien n'établit encore que le gain en plausibilité couvre la perte en validité.

La deuxième est que la \textbf{spécification de la cible est une décision de domaine}, et non un détail de protocole. Le choix retenu en section~\ref{sec:methodo}, la classe standardisée, ne relève pas d'une préférence de notation : il détermine ce qu'un classement de méthodes pourra vouloir dire, puisque deux méthodes évaluées sur des cibles formulées différemment ne sont pas comparables. C'est vraisemblablement l'une des causes du morcellement constaté en section~\ref{sec:etat_art}, et c'est une raison de fixer cette convention avant la campagne plutôt que de la laisser à chaque méthode.

La troisième tient à l'\textbf{ordre de construction}. Commencer par la modélisation mécanistique plutôt que par l'apprentissage profond relevait, sur le moment, de la simple prudence méthodologique : on ne juge pas une performance sans référence. Ce choix a produit un bénéfice qui n'était pas recherché, le second modèle contre lequel s'exerce la cohérence croisée définie en section~\ref{sec:methodo}. Un socle qui serait allé directement à l'apprentissage aurait disposé des mêmes méthodes d'explication et d'aucun moyen de les contrôler.
